\chapter{Metodologia}
\label{cap:metodologia}

Este capítulo descreve a abordagem metodológica que sustenta o trabalho,
as etapas executadas e os critérios adotados em cada uma. A escolha
responde à natureza do problema: não basta entender a teoria do controle
social, nem basta saber construir um Data Warehouse. É preciso fazer os
dois conversarem em um artefato concreto, avaliado contra o problema que
ele se propõe a endereçar.

\section{Abordagem Metodológica: Design Science Research}
\label{sec:dsr}

O trabalho adota a Design Science Research (DSR) como abordagem
metodológica. DSR é orientada à construção e avaliação de artefatos que
respondem a problemas práticos relevantes, articulando rigor científico
e utilidade prática \cite{hevner2004}. Diferentemente de pesquisas
puramente descritivas ou explicativas, a DSR assume que parte do
conhecimento sobre um domínio é produzida \emph{pelo} ato de projetar e
avaliar artefatos nele. Em ciência da computação aplicada, essa
caracterização vem ganhando espaço como alternativa metodológica
rigorosa para trabalhos cuja entrega central é um sistema, um modelo ou
um processo \cite{dresch2015}.

A escolha por DSR é deliberada e responde a duas características deste
trabalho. Primeiro, o produto central é um artefato (o Data Warehouse e
o painel analítico associado), não uma descrição ou uma teoria nova.
Segundo, a relevância do artefato só pode ser aferida contra o problema
que ele endereça, a saber, a distância entre a publicação formal de
dados orçamentários e a possibilidade real de fiscalização cidadã.
Avaliar o artefato é, portanto, parte constitutiva da pesquisa, e não
apêndice opcional.

Hevner e colaboradores \cite{hevner2004} propõem sete diretrizes para
condução de pesquisas em DSR, das quais quatro orientam decisões
concretas deste trabalho:

\begin{itemize}
    \item \textbf{Artefato como pesquisa}: o Data Warehouse e o painel
    analítico são o produto central, documentado no Capítulo
    \ref{cap:dw} em detalhe suficiente para reprodução.

    \item \textbf{Relevância do problema}: a justificativa do trabalho
    parte de lacunas amplamente reconhecidas na literatura de
    transparência pública brasileira (Capítulos \ref{cap:intro} e
    \ref{cap:fundamentacao}), não de curiosidade técnica isolada.

    \item \textbf{Avaliação}: o artefato é avaliado contra
    critérios funcionais explícitos, derivados das perguntas de
    fiscalização levantadas na introdução (Capítulo \ref{cap:avaliacao}).

    \item \textbf{Contribuições claras}: o trabalho declara o que oferece
    de novo em relação ao estado da arte (Capítulo \ref{cap:trabalhos}),
    distinguindo recorte municipal, integração receita-despesa,
    documentação de inconsistências e articulação teórica.
\end{itemize}

O ciclo metodológico segue a estrutura clássica de cinco fases
\cite{dresch2015}: \emph{consciência do problema}, \emph{sugestão},
\emph{desenvolvimento}, \emph{avaliação} e \emph{conclusão}. As fases
não são estanques e o ciclo é iterativo: falhas na fase de
desenvolvimento realimentaram tanto a sugestão quanto a própria
formulação do problema, conforme documentado nas seções a seguir e no
Capítulo \ref{cap:dw}. A correspondência entre fases e capítulos é
direta:

\begin{itemize}
    \item Consciência do problema: Capítulos \ref{cap:intro} e
    \ref{cap:fundamentacao}.
    \item Sugestão: Capítulos \ref{cap:fundamentacao} e
    \ref{cap:trabalhos}.
    \item Desenvolvimento: Capítulo \ref{cap:dw}.
    \item Avaliação: Capítulo \ref{cap:avaliacao}.
    \item Conclusão: Capítulo \ref{cap:conclusao}.
\end{itemize}

\section{Revisão Bibliográfica}
\label{sec:revisao}

A revisão de literatura cobre três frentes que se cruzam no problema:
orçamento público brasileiro (instrumentos, classificações e estágios
da despesa), transparência governamental e controle social, e
modelagem dimensional segundo \citeonline{kimball2013data}.
A literatura sobre civic tech aplicada a dados públicos compõe um quarto
eixo, secundário mas relevante para o posicionamento do trabalho.

A bibliografia cumpre três funções metodológicas distintas. Primeiro,
fornece o vocabulário técnico para descrever o objeto (PPA, LDO, LOA,
estágios da despesa, classificação funcional, MCASP). Segundo,
estabelece a justificativa do problema, com referências empíricas
sobre a distância entre publicação formal e accountability efetiva.
Terceiro, sustenta as decisões de modelagem do artefato, com a
metodologia de Kimball como referência canônica.

A articulação entre literatura técnica e literatura sociopolítica é
escolha deliberada e estabelece a tese central do trabalho: arquitetura
de dados é objeto político. Trabalhos que separam as duas tradições
empobrecem tanto o rigor técnico quanto o impacto cívico.

\section{Pesquisa Documental e Seleção de Fontes}
\label{sec:documental}

A análise documental tem como base primária o Portal da Transparência
da Prefeitura de Juiz de Fora, complementado por duas fontes externas:
o Ementário de Classificação por Natureza da Receita, publicado pela
Secretaria do Tesouro Nacional (STN), e a Lei Orçamentária Anual
publicada pela própria prefeitura.

A seleção das fontes seguiu três critérios. Primeiro, cobertura do ciclo
orçamentário completo: previsão e arrecadação para a receita; empenho,
liquidação e pagamento para a despesa. Segundo, disponibilidade em
formato minimamente processável por pipeline automatizado (planilhas
Excel, mesmo heterogêneas, em vez de PDFs sem estrutura tabular). Quando
formato tabular não estava disponível, como no caso da classificação
funcional da LOA, recorreu-se a OCR sobre PDFs escaneados. Terceiro,
profundidade temporal suficiente para permitir análises de série
histórica.

A pesquisa documental incluiu, ainda, o mapeamento e catálogo dos campos
disponíveis em cada fonte, com documentação explícita das
heterogeneidades encontradas. As inconsistências catalogadas nessa
etapa não foram tratadas como obstáculo a ser silenciosamente
contornado, mas como achado de pesquisa, conforme detalhado na Seção
\ref{sec:inconsistencias}.

\section{Desenvolvimento Iterativo do Artefato}
\label{sec:desenvolvimento}

A fase de desenvolvimento seguiu a metodologia de 
\citeonline{kimball2013data} em quatro passos clássicos: identificação dos
processos de negócio (execução orçamentária de despesas e de receitas),
definição da granularidade dos fatos, identificação das dimensões
relevantes (tempo, unidade administrativa, funcional programática,
função e subfunção LOA, natureza da despesa, natureza da receita,
fornecedor, fonte de recurso e empenho) e identificação dos fatos
(valores empenhado, liquidado e pago para despesa; previsto e
arrecadado para receita).

Esses quatro passos não foram executados em sequência linear. Cada
ciclo de implementação revelou restrições da fonte que obrigaram a
revisitar decisões anteriores. A detecção dinâmica de cabeçalho, o
parser dual para a quebra de layout de 2025-05 e o reconhecimento por
OCR da classificação funcional da LOA, todos documentados no Capítulo
\ref{cap:dw}, nasceram de falhas durante a implementação. Cada falha
foi convertida em decisão de arquitetura documentada, e não em
correção \emph{ad hoc} silenciosa.

Esse comportamento iterativo é coerente com a fase de desenvolvimento
em DSR, que prevê ciclos sucessivos de construção e avaliação parcial
antes da avaliação final do artefato. A escolha da pilha tecnológica
(Dagster para orquestração, Python e pandas para extração, DuckDB como
banco analítico, dbt para modelagem em SQL e Streamlit para os
painéis) responde a critérios de reprodutibilidade local e baixa
barreira de adoção, discutidos na Seção \ref{sec:reprodutibilidade}. As
decisões técnicas detalhadas estão no Capítulo \ref{cap:dw}.

Além das decisões relativas ao \textit{warehouse} propriamente dito, o desenvolvimento incorporou um princípio de design para a camada de consumo que foi privilegiar uma \emph{visão cidadã} sobre uma visão puramente contábil. Na prática, isso se traduziu em três escolhas concretas: agrupar códigos oficiais em categorias mais próximas do vocabulário público (por exemplo, transferências federais e estaduais em lugar de códigos MCASP isolados); acompanhar cada página analítica com um glossário integrado e um bloco orientador com perguntas típicas de fiscalização; e evitar jargão contábil nos títulos e legendas, deslocando-o para \textit{tooltips} acionados por demanda. Essas decisões seguem a leitura freireana desenvolvida no Capítulo~\ref{cap:fundamentacao}, na qual a organização da informação não é neutra e condiciona o universo de leitores possíveis.

\section{Documentação de Inconsistências como Achado}
\label{sec:inconsistencias}

A literatura sobre portais municipais aponta a heterogeneidade dos dados
publicados como obstáculo central à fiscalização
\cite{marco2022transparencia, melo2019proposta}. Este trabalho assume uma
postura metodológica explícita diante desse cenário: inconsistências
encontradas nas fontes são tratadas como achados de pesquisa, não como
ruído a ser corrigido silenciosamente.

Em termos práticos, cada anomalia observada durante a extração foi
registrada e classificada quanto a natureza (mudança de layout, ausência
de padronização de cabeçalho, código não resolvido contra ementário,
divergência entre arquivos do mesmo período) e tratamento aplicado
(parser alternativo, normalização leve, fallback textual, exclusão
documentada). O Capítulo \ref{cap:dw} documenta cada uma dessas decisões
ao lado da seção técnica correspondente.
%, e a Seção \ref{sec:limitacoes} sintetiza as limitações que permanecem.

Essa postura tem fundamento metodológico em DSR: a transparência sobre
as decisões de projeto é parte do que constitui o artefato como
contribuição reprodutível, e não opcional cosmético. Tem também
fundamento normativo: trabalhos sobre transparência pública que
silenciam sobre suas próprias decisões de tratamento reproduzem
exatamente o problema que pretendem combater.

\section{Verificação por Invariantes de Domínio}
\label{sec:invariantes}

Em dados fiscais publicados, falha silenciosa custa caro: um valor
negativo por erro de parsing, uma quebra de monotonicidade no acumulado
anual ou uma dimensão com chave duplicada pode virar conclusão errada
sobre execução orçamentária. A verificação da qualidade dos dados foi,
portanto, incorporada como etapa metodológica explícita, e não como
detalhe técnico.

A verificação opera em duas camadas. A primeira aplica testes
estruturais sobre o modelo dimensional (unicidade e não-nulidade de
chaves, integridade referencial entre fatos e dimensões, presença das
medidas obrigatórias), declarados em arquivos de configuração do dbt e
executados a cada materialização. A segunda aplica invariantes de
domínio expressas em SQL: a arrecadação acumulada no exercício deve ser
não-decrescente mês a mês; o valor a realizar deve coincidir com a
previsão atualizada menos o arrecadado; o acumulado anual não pode ser
inferior ao acumulado dos meses anteriores na mesma natureza de receita.

A escolha dessas invariantes não é arbitrária: cada uma corresponde a uma
classe específica de falha plausível dada a forma como o portal publica
os dados (carga parcial entre meses, parsing equivocado, drift de
schema). O catálogo completo dos testes está no Capítulo \ref{cap:dw}.

\section{Avaliação Funcional Comparativa}
\label{sec:metodo_avaliacao}

A fase de avaliação do ciclo DSR foi implementada como uma análise
funcional comparativa entre o Portal da Transparência da Prefeitura de
Juiz de Fora e o painel analítico desenvolvido. O critério adotado é
funcional, e não estético: dada uma pergunta de fiscalização, quantos
passos cada ambiente exige para respondê-la, com que grau de
consolidação manual e com que nível de mediação técnica intermediária?

\subsection*{Seleção de tarefas}

As tarefas avaliadas foram derivadas das perguntas exemplares levantadas
na introdução do trabalho e da literatura sobre uso de portais por organizações de fiscalização \cite{marco2022transparencia}. O conjunto cobre quatro classes de operação:

\begin{itemize}
    \item Agregação simples sobre uma única dimensão (por exemplo,
    pagamentos totais por função de governo em um exercício).
    \item Comparação temporal entre exercícios (por exemplo, evolução de
    investimento em uma área ao longo de três anos).
    \item Cruzamento entre fontes (por exemplo, previsão versus
    arrecadação por natureza de receita).
    \item Análise de concentração (por exemplo, participação dos cinco
    maiores fornecedores no total pago).
\end{itemize}

A escolha das classes assegura que a avaliação não fique restrita a
operações em que qualquer ferramenta analítica trivialmente supera uma
planilha bruta.

\subsection*{Critérios de comparação}

Cada tarefa é avaliada em três dimensões:

\begin{itemize}
    \item \textbf{Número de passos manuais} necessários no portal versus
    no painel.
    \item \textbf{Necessidade de consolidação} entre arquivos distintos
    do portal.
    \item \textbf{Grau de mediação técnica} requerido (por exemplo,
    interpretação de códigos MCASP, decodificação de funcional
    programática).
\end{itemize}

O resultado consolidado da comparação está na Tabela
\ref{tab:comparacao-portal-dashboard} do Capítulo \ref{cap:avaliacao}.

\subsection*{Limitações da avaliação}

A avaliação é autoral. Não foram realizados testes de usabilidade com
usuários finais, e a comparação foi conduzida pelo próprio autor com
base na reprodução repetida das tarefas em ambos os ambientes. Essa
limitação é reconhecida explicitamente: a evidência produzida diz
respeito à \emph{redução de fricção operacional} entre portal e painel,
não à compreensão efetiva dos dados por cidadãos sem formação
especializada.

A opção pela avaliação autoral tem três justificativas dentro do escopo
deste trabalho. Primeiro, o protocolo de comparação é reprodutível: as
tarefas e os critérios estão documentados e podem ser reaplicados por
terceiros. Segundo, as tarefas escolhidas estão ancoradas na literatura
sobre fiscalização cidadã, e não em preferência subjetiva. Terceiro,
um estudo com usuários exigiria submissão a comitê de ética e
recrutamento estruturado, fora do escopo de tempo de um TCC. Essa
extensão fica registrada como trabalho futuro no Capítulo
\ref{cap:conclusao}.

\section{Compromisso com Reprodutibilidade}
\label{sec:reprodutibilidade}

Reprodutibilidade foi assumida como princípio metodológico desde o início do projeto, e não como entrega adicional. Para um trabalho que se posiciona como referência metodológica replicável para outros municípios, reprodutibilidade não é opcional, mas sim parte da contribuição.

O compromisso se operacionaliza em quatro decisões concretas. Primeiro, todo o código está versionado em repositório público, com instruções de execução e descrição da pilha tecnológica. Segundo, o pipeline foi projetado para ser idempotente: reexecutar uma partição produz o mesmo estado final do warehouse, sem duplicatas nem lacunas. Terceiro, as chaves substitutas das dimensões são geradas por hash determinístico sobre campos naturais, de modo que o reprocessamento preserva identidade entre execuções. Quarto, o painel analítico está publicado em instância pública (Streamlit Community Cloud), permitindo que qualquer leitor verifique os resultados sem instalar o ambiente localmente.
Os mecanismos técnicos correspondentes estão detalhados no Capítulo \ref{cap:dw}.

\section{Estudo de Caso: Juiz de Fora}
\label{sec:estudo_caso}

A aplicação do modelo a Juiz de Fora cumpre dois papéis dentro do ciclo
DSR. Primeiro, valida o modelo dimensional proposto contra um caso
concreto: se o esquema responde às perguntas de fiscalização derivadas
da literatura, ele serve. Segundo, produz documentação dos desafios
técnicos encontrados em um município específico, oferecendo referência
para futuras adaptações em outras cidades.

A escolha de Juiz de Fora não é casual e responde a três fatores: a
vinculação institucional do trabalho (UFJF); a existência de portal
municipal ativo com dados publicados em formatos minimamente
analisáveis; e o porte intermediário do município, que torna o caso
representativo para um universo maior de cidades brasileiras de porte
médio com infraestrutura digital comparável.
